\renewcommand{\chapcolor}{chMesh}
\section{Service Mesh}
\subsection{Microservices}
Small independent services, one business function, communicate via lightweight APIs. Microservices == distributed computing. \\
\textbf{Pros:} independent deploy, flexibility, scalability, fault isolation, team autonomy. \\
\textbf{Cons:} operational complexity, distributed-system challenges, data consistency, higher resource use. \\
\textbf{Challenges:} \textit{Communication} - services must know endpoints, new service ripples changes. \textit{Security} - in-cluster traffic unsecured, any pod $\to$ any pod. \textit{Monitoring} - spread out, non-business logic added per app.

\subsection{Distributed Computing Fallacies}
Network reliable; latency zero; bandwidth infinite; topology stable; one admin; transport cost zero; network homogeneous. All false $\to$ motivate mesh.

\subsection{Service Mesh}
Dedicated infra layer handling service-to-service communication: traffic routing, security, observability, resiliency - abstracted from app code. \\
\textbf{Capabilities:} Traffic Mgmt (routing, splitting), Resilience (fault tolerance, retries, fallback), Security (mTLS, authN/Z), Observability (metrics, logs, tracing). \\
\textbf{Why:} centralized policy, uniform observability, faster incident resolution, transparent security, advanced rollouts (canary/AB), resilience.

\textbf{Components:} \textit{Data Plane} - lightweight proxies (Envoy), handle comm/metrics/security. \textit{Control Plane} - manages config/policies (Istiod), pushes rules to data plane.

\subsection{Mesh Evolution: Library $\to$ Sidecar $\to$ Kernel}
\textbf{Library:} mesh logic linked into each app, per-language lib (Python, Go, ...). Invasive, language-locked. \\
\textbf{Sidecar:} proxy per pod, transparent traffic interception, no code change, language-agnostic. Downsides: invasive pod-spec mutation, upgrade = pod restart, large CPU/RAM per pod (worst-case provisioning), extra hop. \\
\textbf{Kernel/Ambient:} mesh in kernel (eBPF) or per-node agent. 30 pods/node $\to$ 1 proxy/node instead of 30. L7 still typically needs userspace proxy.

\textbf{Leading meshes:} \textit{Istio} (Google, feature-rich), \textit{Linkerd} (lightweight, perf), \textit{Kuma} (CNCF, VMs+containers), \textit{Consul Connect} (HashiCorp, discovery), \textit{Cilium} (eBPF, not dedicated mesh).

\subsection{Sidecar Pattern Pros/Cons}
\textbf{Pros:} + transparent (no code change, lang-agnostic), + per-pod isolation, + platform team owns cross-cutting (mTLS/retries/routing/telemetry). \\
\textbf{Cons:} - CPU/RAM scales with pod count, - extra hop = latency, - sidecar injection + proxy upgrade needs restart, - overkill for few services.

\includegraphics[width=0.19\textwidth]{ServiceMesh.png}

Envoy proxies authenticate via \textbf{mTLS} (mutual TLS, both sides present certs).

\subsection{Ambient Mesh (Sidecarless Istio)}
Sidecar downsides $\to$ ambient: per-node \textbf{zTunnel} agent (L4, mTLS, authN/Z, no HTTP parsing) communicates across nodes via HTTPS CONNECT tunnel. Optional L7 \textbf{Waypoint Proxy} per namespace (Envoy) for rich policies/telemetry. zTunnel alone = L4 metrics only.

\textbf{Slicing:} \textit{Secure Overlay (L4):} TCP routing, mTLS tunneling, simple authZ, TCP metrics. \textit{L7 Processing:} HTTP routing/LB, circuit breaking, rate limiting, fault injection, retry, timeouts, rich authZ, HTTP metrics/access logs/tracing.

\textbf{Cilium variant:} dynamic \textbf{eBPF} programs in kernel handle L3/L4 (network path); per-node Envoy proxy for L7 termination when needed. No sidecar, no extra hop for L4.

\subsection{Gateway API \& Ambient Routing / GAMMA}
Enable ambient: \texttt{kubectl label ns x istio.io/dataplane-mode=ambient} (sidecar mode: \texttt{istio-injection=enabled}). Routes via K8s \textbf{Gateway API} (\texttt{HTTPRoute}) not VirtualService. Two Deployments behind one Service $\approx$50/50 default; override via weights. Best practice: \texttt{version} label per Deployment. \\
\textbf{GAMMA} (Gateway API for Mesh Management \& Administration): Gateway API for east/west (mesh) traffic, not just north/south (ingress). Adopted by more meshes.
\begin{lstlisting}[language=yaml]
# weighted canary 90/10:
- backendRefs:
  - { name: frontend,    port: 5000, weight: 90 }
  - { name: frontend-v2, port: 5000, weight: 10 }
# mirror/shadow (copy to v2, drop response):
- backendRefs:
  - { name: frontend, port: 5000 }
  filters:
  - type: RequestMirror
    requestMirror: { backendRef: { name: frontend-v2, port: 5000 } }
\end{lstlisting}
\textbf{istio-system pods:} Istiod, Ingress Gateway (Envoy), Prometheus, Grafana, Jaeger, Kiali. \textbf{Kiali shapes:} rectangle=app, triangle=Service, square=workload/pod.

\subsection{Traffic Management (Istio Classic)}
Gateway replaces K8s ingress/egress. VirtualService binds gateway, routes to services. DestinationRule applies post-routing policy: LB, connection pool, outlier detection (circuit breaking), TLS mode, defines subsets (versions).
\begin{lstlisting}[language=yaml]
apiVersion: networking.istio.io/v1beta1
kind: Gateway
metadata: { name: my-gateway }
spec:
  selector: { istio: ingressgateway }
  servers:
  - port: { number: 80, name: http, protocol: HTTP }
    hosts: ["example.com"]
--
kind: VirtualService
metadata: { name: my-vs }
spec:
  hosts: ["example.ch"]   # must match gateway host
  gateways: [my-gateway]
  http:
  - route:
    - destination:
        host: my-service
        weight: 99        # split traffic if 2nd dest defined
        port: { number: 80 }
--
kind: DestinationRule
metadata: { name: my-dr }
spec:
  host: my-service
  trafficPolicy:
    tls: { mode: ISTIO_MUTUAL }
\end{lstlisting}

\subsection{Resilience \& Use Cases}
\textbf{Fault Injection:} simulate net failures, test error handling. \textbf{Timeouts:} cap request time, avoid queue buildup. \textbf{Retries:} auto-retry on failure. \\
\textbf{Cascading Failure:} one service dies $\to$ callers pile up $\to$ whole graph red. \\
\textbf{Circuit Breaking:} trip on errors/outlier $\to$ prevent cascading failure (DestinationRule outlierDetection). \\
\textbf{Rate Limiting:} cap RPS, protect downstream from overload. \\
\textbf{Zero-Trust:} encrypt + authenticate every call, no implicit trust in cluster. \\
Istio exports Prometheus metrics by default.

\subsection{Deployment Strategies}
\textbf{Rolling:} gradual pod replace, default K8s. \textbf{Recreate / Big-bang:} kill all, deploy new (downtime). \textbf{Canary:} small \% to v2, watch, ramp up. \textbf{Blue/Green:} two envs, switch traffic atomically, easy rollback. \textbf{Dark Launching / Shadow:} mirror prod traffic to v2, drop response (test under real load). See \textit{General Definitions}.

\subsection{Security}
\textbf{AuthN:} mTLS between services. \textbf{AuthZ:} AuthorizationPolicy = fine-grained RBAC.
\begin{lstlisting}[language=yaml]
# enforce mTLS in namespace
kind: PeerAuthentication
metadata: { name: default, namespace: production }
spec:
  mtls: { mode: STRICT }
--
kind: AuthorizationPolicy
metadata: { name: require-jwt, namespace: foo }
spec:
  selector: { matchLabels: { app: httpbin } }
  action: ALLOW
  rules:
  - from: [{ source: { namespaces: ["frontend"] } }]
    to:   [{ operation: { methods: ["GET"], paths: ["/headers"] } }]
\end{lstlisting}
